\documentclass[10pt,a4paper]{amsart}
\usepackage[margin=19mm]{geometry}
\usepackage{amsmath,amssymb,mathtools}
\usepackage[scr=rsfs,cal=euler]{mathalfa}
\usepackage{xfrac}
\usepackage[unicode,hidelinks,bookmarks=false]{hyperref}
\newcommand{\ie}{i.e.\ }
\newcommand{\cf}{cf.\ }
\newcommand{\resp}{respectively\ }
\newcommand{\Subsection}{\subsection}
\providecommand{\nobreakdash}{\nobreak\hbox{-}}
\setlength{\emergencystretch}{2em}
\multlinegap0pt
\usepackage[shortlabels]{enumitem}
\usepackage{amssymb}
\usepackage{xcolor}
\usepackage{float}
\usepackage{tikz}
\usepackage{mathtools}
\newtheorem{theorem}{Theorem}
\newcommand{\K}{\mathbb{K}}
\newcommand{\N}{\mathcal{N}}
\newcommand{\map}{\varrho}
\newcommand{\pp}{\prec\hspace{-2pt}\prec}
\newcommand{\prob}{\pi}
\newcommand{\rw}{\beta}
\newcommand{\subsV}{U}
\title{Recurrence and transience for non-Archimedean and directed graphs}
\author{Matthias Keller, Anna Muranova}
\date{Source: arXiv:2406.17344v1 (CC BY 4.0)}
\begin{document}
\maketitle

\noindent\emph{Source and licence.} Recurrence and transience for non-Archimedean and directed graphs. Version arXiv:2406.17344v1 (CC BY 4.0), distributed by arXiv under the Creative Commons Attribution 4.0 International licence, http://creativecommons.org/licenses/by/4.0/, as recorded in the arXiv licence metadata for this exact version and shown on the abstract page https://arxiv.org/abs/2406.17344v1. Full licence notice: \texttt{licenses/arxiv-2406.17344-CC-BY-4.0.txt}.\par\smallskip

\noindent\textit{Editorial scope.} Counted unit: the article's theorem thm:finitevsdirected (source0.tex lines 798--805). The article's complete original proof of it is delivered with it (source0.tex lines 807--860, one \texttt{proof} environment of 54 lines). No mathematics is written, repaired or re-derived: the statement and the proof are the article's own lines, in the article's own notation, and no retained line is substituted or re-flowed.  Retained uncounted as the source-local context the proof needs: lines 590--613.  Nothing was omitted from any retained span, so the package carries no omission record.  No retained line contains a citation, so the wrapper carries no bibliography and no bibliography entry.  No retained line contains a figure environment, an image-inclusion command or drawing code, so no image is delivered and the package declares no asset dependency.  Editorial formatting only, in addition: an AMS \texttt{amsart} preamble that restates the article's own declarations and macros used by the retained lines, namely the environments \texttt{theorem} on separate counters, together with the standard packages those lines need.  The retained environments print in this document as Theorem 1 in order of appearance, whereas the article prints the same items with the numbering of its own full document.  The complete native source of the article as distributed in the arXiv e-print for this exact version is bundled unchanged as \texttt{sources/161328/source0.tex}.
\begin{theorem}\label{thm::bxby}
Let $b$ be a  graph over $V$, $\subsV\subseteq V$, $x,y\in \subsV$ and $\pi=\pi_\subsV^{b}$. Then the following holds:
\begin{enumerate}[label=${(\alph*)}$]
\item We have 
$x\to y$ if and only if there exists a path $$ x=x_0\sim x_1\sim\dots \sim x_n=y $$ in  $\subsV$ (with respect to $ b $) such that for all $ i=0,\ldots,n-1 $
 $$ b(x_i,x_{i+1})\simeq b(x_i). $$
In particular, in this case 
${b(x_{i+1})\succsim b(x_i)}$, $ i=0,\ldots,n-1$
and ${b(y)}\succsim{b(x)}$.
Furthermore, in this case this path is also a directed path (with respect to $ \pi $)
$$
x=x_0\leadsto x_1\leadsto\dots \leadsto x_n=y
$$
and any such directed path is a path with respect to $ b $ with the properties above. 
\item
Let $x\to y$. Then $y\to x$ if and only if  $b(x)\simeq b(y).$


\item
 We have  $x\leftrightarrow y$ if and only if $b(x)\simeq b(y)$ and there exists a path  $$ x=x_0\sim x_1\sim\dots \sim x_n=y $$ in  $\subsV$  such that for all $ i=0,\ldots,n-1 $
with $$ b(x_i,x_{i+1})\simeq b(x). $$ 
In particular, this path is also a directed path form  $ x $ to $ y $ and from $ y $ to $ x $.
\end{enumerate}
\end{theorem}

\begin{theorem}\label{thm:finitevsdirected} Let $ \pi $ be a transition matrix over a finite set $ V $. Then, there exists a weighted graph $ b $ over $ V $ with respect to a non-Archimedean ordered field $\Bbb K$ and $\subsV\subseteq V$ such that $\pi=\pi_{U}^{b}$ if and only if  there is a function $\rw:V\to \Bbb R^+$ such that 
	for any irreducible class $C$ and  $x,y\in C$ 
	$$
	\prob(x,y)\rw(x)=\prob(y,x)\rw(y)
	$$
	and $ \pi(x,x)=0 $ if $ x $ is not an absorbing state.
	Furthermore, for a given transition matrix, $ \beta $ can be chosen such that $ \beta (a)=1 $ for some $ a\in C $ and $  \beta(x)=\rho (b(x)/b(a)) $ for all $ x\in C $.
\end{theorem}

\begin{proof}
	We  construct $b$ and $U$. First, we fix an infinitely large $\N \in\K$, i.e., $1\pp\N$. Furthermore, we know that the set of irreducible classes is a finite partially ordered set and has a height function $ h $.  For  $ x \in  V$ contained in a irreducible class  $ C_{x} $, we denote $ h(x)=h(C) $.
	We let $ U $
	be the subset of $ V $ which are not absorbing states. To define $ b $,  we let $ b=0 $ on $ (V\setminus \subsV) \times(V\setminus \subsV) $. Furthermore, we define
	for  $x\in U$,  $y\in C_{x}$
	\begin{align*}
		b(x,y) 
		=	\prob(x,y)\rw(x) \N^{h({x})}
	\end{align*}
and  for $ y\in V\setminus C_{x} $
	\begin{align*}
b(x,y) 
=		\max\{\prob(x,y)\rw(x),\prob(y,x)\rw(y)\}\N^{\min\{h({x}), h({y})\}}.
\end{align*}
Then $ b $ is clearly symmetric and has zero diagonal.
Observe,  for $ x\in U $,
\begin{align*}
	b(x,y)=
	\begin{cases}
		\prob(x,y)\rw(x) \N^{h({x})} &: x\leadsto y, \mbox{ i.e., $ \prob(x,y)>0 $}\\
		\prob(y,x)\rw(y) \N^{h({y})} &: \mbox{else}.
	\end{cases}
\end{align*}
Note that in the second case we have $ b(x,y)=0 $ whenever $ y\not\leadsto x $ and $ h(y)<h(x) $ whenever $ y\leadsto x $.
As $ \prob(x,x)=0 $ for $ x\in U $, we have 
\begin{align*}
	  b(x)  &=\sum_{y:x\leadsto y}b(x,y)+\sum_{\substack{y:x\not\leadsto y,\\ \;\; y\leadsto x }}b(x,y)=\sum_{ y\in V } 	\prob(x,y)\rw(x) \N^{h({x})} +
\sum_{\substack{y:x\not\leadsto y,\\ \;\; y\leadsto x}} 	\prob(y,x)\rw(y) \N^{h({y})}\\
& =  \left(  1+ \sum_{\substack{y:x\not\leadsto y,\\ \;\; y\leadsto x} } 	\prob(y,x)\frac{\rw(y)}{\rw(x)} \N^{- (h(x)-h({y}))}\right)\rw(x)\N^{h(x)}=(1+\tau)\rw(x)\N^{h(x)}
\end{align*}
for some infinitesimal $ \tau $. We obtain for $ \prob(x,y)>0 $ and $ x\neq y $ (which implies $ x\in U $)
\begin{align*}
	\pi^{b}_{\subsV}(x,y) = \rho(p(x,y))=\rho\left(\frac{b(x,y)}{b(x)}\right)
	=\rho\left(\frac{\prob(x,y)\beta(x)\N^{h(x)}}{(1+\tau)\beta(x)\N^{h(x)}}\right)
	=\rho\left(\frac{\prob(x,y)}{(1+\tau)}\right)=\prob(x,y)
\end{align*}
and for $ \prob(x,y)=0 $ (which includes the case $ x\in V\setminus U $ and $ x\neq y $) where  we have either $ \prob(y,x)=0 $ or $ h(y)<h(x) $
\begin{align*}
	\pi^{b}_{\subsV}(x,y) =\rho\left(\frac{\prob(y,x)\beta(y)\N^{h(y)}}{(1+\tau)\beta(x)\N^{h(x)}}\right)= \rho\left(N^{-(h(x)-h(y))}\frac{\prob(y,x)\beta(y)}{(1+\tau)\beta(x)}\right)=0=\prob(x,y).
\end{align*}
Finally, $ \pi_{U}^{b}(x,x) =1$ by definition for $ x\in V\setminus \subsV $.\medskip




To prove the opposite direction we fix in each  irreducible class $C$ a vertex $x\in C$ and put
 $\rw(x)=1.
$
Further, for any $y\in C$, we prescribe $\rw(y)=\map({b(y)}/{b(x)})$. Then, we have, for any $y,z\in~C$,
\begin{align*}
\prob(y,z) \rw(y)=\map(p(y,z))\map\left(\dfrac{b(y)}{b(x)}\right)=\map\left(\dfrac{p(y,z) b(y)}{b(x)}\right)=\map\left(\dfrac{p(z,y) b(z)}{b(x)}\right)=\prob(z,y) \rw(z),
\end{align*}
where the second and last equality hold, since $p\preceq 1$ and ${b(y)}/{b(x)}$ as well as ${b(z)}/{b(x)}$  have no infinitely large component, since $y,z$  and $ x$ are all in the same irreducible class  $C$, confer Theorem~\ref{thm::bxby}.
\end{proof}

\end{document}
